\section{路径 A：端到端评估扩容}

\subsection{A0：恢复执行已发布的 \texorpdfstring{$n{=}100$}{n=100} 协议（P0，零新工程）}

第一步不写任何新代码。直接挂起限流感知驱动器：\texttt{python3 scripts/run\_qa\_driver.py 540 760}（窗口 540 秒、目标 760 次成功调用），断点文件 \texttt{qa\_expanded.jsonl} 会自动接续已有进度。按头部文档记载的调用预算：总计 700 个方法--记录对（100 记录 $\times$ 7 方法），每问最多 96 个答案 token 加提示本身（中位上下文约 5{,}000 词，压缩后 $4\times$ 约 1{,}250 词），合计输入 token 量级约 $700\times1{,}500\approx 10^6$、输出 $700\times 96\approx 6.7\times 10^4$。以当前生产 API 的计费量级，这是一次可一次性完成的预算；执行节奏由驱动器按限流自适应，建议在配额宽松时段连续挂机，断点机制保证中断无损失。

A0 完成后立即产出三件可写进修订稿的东西：其一，$n{=}100$ 的端到端主表（七方法 $\times$ EM/F1/保留率，附配对 bootstrap 区间），取代 $n{=}25$ 版本的示踪性区间；其二，BM25 首次进入端到端比较（$n{=}25$ 版没有它），``结构性权衡''叙事由此获得端到端层的对应物——若 BM25 端到端与召回层同样领先，权衡叙事在两层一致；若不领先，则获得``召回优势不传导到答案质量''的第二类发现，同样有价值；其三，逐记录召回--F1 配对扩容到 $n{=}100$（与门同集），相关系数的区间收窄约一半。

\subsection{A1：扩容到 \texorpdfstring{$n{=}300$}{n=300}（P1）}

样本量论证：$n{=}25$ 时 \sieve{}--截断配对差的区间半宽约 $\pm 11.4$ F1 点；配对区间半宽按 $1/\sqrt{n}$ 收缩，$n{=}100$ 时约 $\pm 5.7$，$n{=}300$ 时约 $\pm 3.3$。若端到端效应真实存在于 $\pm 5$ 点的量级（召回层配对增益 12.8 点的合理传导折扣），$n{=}300$ 是能以 95\% 置信度分离它的最低规模；若区间仍然跨零，``统计持平''的主张便有了与其措辞相称的证据，同样构成完成。实现上只需把 \texttt{run\_qa\_expanded.py} 的测试集从``前 100 条''参数化为``前 300 条''（池本身有 1{,}000 条，无需重建），七方法不变、$r{=}4$ 不变、提示与解码协议逐字不变——协议冻结是本路径的纪律：任何对提示模板、温度、答案长度约束的改动都会使新旧数据不可合并。预算为 A0 的三倍量级，仍属单次可完成。

\subsection{A2：跨数据集端到端（P2，视资源）}

在 HotpotQA 与 2WikiMultihopQA 池上各抽 $n{=}100$（各池均有 300 条，抽样种子固定并登记）。这两池的上下文是多文档拼接而非单篇论文，端到端行为可能与 QASPER 不同——多跳问题在 $4\times$ 压缩下可能更脆弱，而这正是 ``problem formulation'' 声明（词汇信号强度决定问题条件化的收益）在端到端层的直接检验点。方法集可精简为四种（full/head/bm25/sieve）以控制预算：full 与 head 定义端到端上下界，bm25 与 sieve 是权衡叙事的两个端点，random/stride/textrank 在 $n{=}100$ 下的区分度不足，省去。预算约为 A0 的两倍（$2\times100\times4=800$ 对）。

\subsection{A3：结果表格模板（先于实验固定）}

端到端主表在跑数之前先定格式并冻结列含义：行为方法（full context、head、random、stride、textrank、bm25、sieve），列为 EM、F1、full-context F1 保留率、输入 token 中位数；区间一律配对 bootstrap（1{,}000 重采样、95\%），\sieve{} 行与每个基线行的差值区间单独成列或成表。跨数据集扩展沿用同表式样。表格模板先行的作用是把``要报告什么''与``跑什么''解耦，避免事后挑数。
